\section{Recursive forcing and the truth lemma}
\label{sec:p06-forcing-recursion-truth}

The recursive relation is constructed before it is compared with truth in an
extension. Its atomic definition uses name ranks, not generic filters or a
truth predicate. The truth lemma then supplies the set bounds needed in the
Separation, Replacement and Power Set arguments. This is a modern
reconstruction of the density-closed presentation in
\cite[Sections~2.1--2.2, pp.~6--10]{Karagila2026ForcingNotes}, with the
Collection and subset-name arguments compared with
\cite[Sections~10.1--10.3, pp.~26--36]{Unger2014ForcingSummerSchool}.
It is not an identification of either text with the separately cited
\cite{Jech2003,Kunen2011} presentations.

\begin{hbconvention}{def:p06-ftr-ground}{Ground domain and supplied generic filter}
Work in classical \(\mathsf{ZFC}\) as metatheory. Let \(M\) be a transitive
set such that the membership structure \(\mathcal M=(M,\in)\) satisfies
\(\mathsf{ZFC}\). Let \(\mathbf P=(P,\leq_P,1_P)\in M\) be a nonempty
partial order with greatest element; smaller conditions are stronger.
No separativity, chain condition, closure, or external countability is assumed.
Compatibility means existence of a common strengthening in \(P\).
The poset-filter and dense-set clauses are those of
\cref{def:p06-poset-filter}. A filter \(G\subseteq P\) is supplied in the
ambient universe and meets every dense subset of \(P\) belonging to \(M\).
Its existence is an additional hypothesis, not a conclusion from transitivity.

All name and forcing constructions carrying the subscript \(M\) are computed
in \(\mathcal M\). Quantification over names in such a clause ranges over
names in \(M\), not over additional names of the ambient universe. The
structure \((M[G],\in)\) is evaluated externally, formula by formula.
\end{hbconvention}

\subsection{Set bounds and name ranks}

\begin{hbdefinition}{def:p06-ftr-names}{Names, their closure and intrinsic rank}
Retain the name--condition order of the earlier name definition:
an entry in \(\sigma\) is \(\langle\pi,r\rangle\), where \(\pi\) is a
subname and \(r\in P\). Write
\[
 A_\sigma=\{\pi:\exists r\in P\ (\langle\pi,r\rangle\in\sigma)\}.
\]
Close \(\{\sigma\}\) under these first-coordinate subnames by an
\(\omega\)-sequence of unions; call the resulting set \(S_\sigma\).
The predicate \(\mathsf{Name}_P(\sigma)\) says that every member of every
\(\tau\in S_\sigma\) is an ordered pair whose second coordinate is in
\(P\). The construction of the closure rejects malformed entries.
Its equivalent hierarchy definition uses the earlier stages, computed in
the ground structure:
\[
 W_\alpha^P=
 \mathcal P^M\!\left(\left(\bigcup_{\beta<\alpha}W_\beta^P\right)\times P\right).
\]
Here \(\mathcal P^M\) is the power set supplied by \(\mathcal M\), and
\(W_0^P=\{\varnothing\}\). This is the earlier all-ordinal stage convention,
not a silent replacement by a continuous hierarchy with empty stage zero.

The intrinsic name rank is defined by
\[
 \rho_P(\sigma)=\sup\{\rho_P(\pi)+1:\pi\in A_\sigma\}.
\]
Empty suprema are zero. Evaluation and canonical names retain
\[
 \sigma^G=\{\pi^G:\exists r\in G\ (\langle\pi,r\rangle\in\sigma)\},
 \qquad \check x=\{\langle\check y,1_P\rangle:y\in x\}.
\]
These are recursion specifications; the next result establishes their
existence and exact set bounds.
\end{hbdefinition}

\begin{hbtheorem}{thm:p06-ftr-name-foundations}{proved=annex:p06.ftr.name-foundations}{Names and evaluation are well defined}
The name predicate is first-order definable. Its closure test and hierarchy
definition agree. Every ground name has a unique rank in \(M\), and
\[
 \sigma\in W_\alpha^P\quad\Longleftrightarrow\quad\rho_P(\sigma)\leq\alpha
 \quad\text{inside }\mathcal M.
\]
Ranks and the property of being a name agree with the ambient computation
on names in \(M\). Evaluation has a unique value for every ground name.
For any poset filter containing \(1_P\),
\(M[G]=\{\sigma^G:\sigma\in M,\ \mathsf{Name}_P(\sigma)\}\) is a
transitive set, contains \(M\) and \(G\), and
\[
 \operatorname{rank}(\sigma^G)\leq\rho_P(\sigma).
\]
Consequently \(M[G]\) and \(M\) have the same ordinals. These conclusions
do not yet assert that an arbitrary nongeneric filter preserves
\(\mathsf{ZFC}\).
\end{hbtheorem}

\subsection{The atomic recursion and all logical constructors}

\begin{hbdefinition}{def:p06-ftr-atomic}{The density-closed atomic clauses}
Write \(p\forces^{\mathrm{rec}}_{M,P}\varphi\) for the new relation.
The superscript distinguishes it from the earlier external semantic
relation. In the following clauses all names and conditions belong to
the ground structure. Define the auxiliary inclusion test by
\[
 \mathsf{Incl}_P(p,\sigma,\tau)
 \quad\Longleftrightarrow\quad
 \forall\langle\pi,r\rangle\in\sigma\ \forall q\leq_P p,r\
       \bigl(q\forces^{\mathrm{rec}}_{M,P}\pi\in\tau\bigr).
\]
It is an auxiliary test on names, not an extra primitive of the membership
language. The two atomic clauses are
\[
\begin{aligned}
 p\forces^{\mathrm{rec}}_{M,P}\sigma=\tau
 &\Longleftrightarrow
   \mathsf{Incl}_P(p,\sigma,\tau)\land\mathsf{Incl}_P(p,\tau,\sigma),\\
 p\forces^{\mathrm{rec}}_{M,P}\sigma\in\tau
 &\Longleftrightarrow
   \{q\leq_P p:\exists\langle\pi,r\rangle\in\tau\
      (q\leq_P r\land q\forces^{\mathrm{rec}}_{M,P}\sigma=\pi)\}
       \text{ is dense below }p.
\end{aligned}
\]
The recursion measure is the ordered pair
\[
 \bigl(\max(\rho_P(\sigma),\rho_P(\tau)),
       \min(\rho_P(\sigma),\rho_P(\tau))\bigr)
\]
in lexicographic order. Replacing either name by a proper subname strictly
decreases this measure, even when the other name has larger rank.
\end{hbdefinition}

\begin{hbdefinition}{def:p06-ftr-logical}{Logical recursion and finite formula scope}
For each fixed standard membership-language formula, first expand
\(\varphi\lor\psi\) as \(\neg(\neg\varphi\land\neg\psi)\),
\(\varphi\to\psi\) as \(\neg(\varphi\land\neg\psi)\), and
\(\forall x\,\varphi\) as \(\neg\exists x\neg\varphi\), avoiding capture.
Use these primitive clauses:
\[
\begin{aligned}
 p\forces^{\mathrm{rec}}_{M,P}\bot&\quad\text{never},\\
 p\forces^{\mathrm{rec}}_{M,P}\neg\varphi
 &\Longleftrightarrow\neg\exists q\leq_P p\
                      (q\forces^{\mathrm{rec}}_{M,P}\varphi),\\
 p\forces^{\mathrm{rec}}_{M,P}\varphi\land\psi
 &\Longleftrightarrow(p\forces^{\mathrm{rec}}_{M,P}\varphi)
                     \land(p\forces^{\mathrm{rec}}_{M,P}\psi),\\
 p\forces^{\mathrm{rec}}_{M,P}\exists x\,\varphi(x,\bar\tau)
 &\Longleftrightarrow
 \{q\leq_P p:\exists\sigma\
   (\mathsf{Name}_P(\sigma)\land
      q\forces^{\mathrm{rec}}_{M,P}\varphi(\sigma,\bar\tau))\}
      \text{ is dense below }p.
\end{aligned}
\]
The existential name quantifier is an ordinary ground-model set quantifier.
It is not a quantifier over generic filters or over class objects.
Each fixed formula has its own defining formula
\(\mathsf{Frc}_{P,\varphi}(p,\bar\tau)\); no internal predicate is claimed
to work uniformly for all formula codes.
\end{hbdefinition}

\begin{hbtheorem}{thm:p06-ftr-recursion}{proved=annex:p06.ftr.recursion}{Existence, uniqueness and definability of recursive forcing}
The atomic recursion determines unique truth sets of conditions, compatible
on overlapping rank-bounded name sets. It is defined by a ground-model
first-order formula. For each fixed standard formula the finite logical
recursion then supplies \(\mathsf{Frc}_{P,\varphi}\), using only the ground
poset, names and the stated set-theoretic operations. No truth lemma, generic
existence or satisfaction predicate for the universe is used in this
construction.
\end{hbtheorem}

\begin{hbtheorem}{thm:p06-ftr-density}{proved=annex:p06.ftr.density}{Monotonicity, density closure and decision}
For every fixed formula and tuple of ground names:
\begin{enumerate}
\item If \(p\) forces it and \(q\leq_P p\), then \(q\) forces it.
\item \(p\) forces it exactly when the conditions forcing it are dense below \(p\).
\item The set of conditions forcing it or its negation is dense and belongs
      to \(M\). No condition forces both.
\end{enumerate}
The expanded constructors have these exact clauses:
\[
\begin{aligned}
 p\forces^{\mathrm{rec}}_{M,P}\varphi\lor\psi
 &\Longleftrightarrow
   \{q\leq_P p:q\forces^{\mathrm{rec}}_{M,P}\varphi
                  \text{ or }q\forces^{\mathrm{rec}}_{M,P}\psi\}
                 \text{ is dense below }p,\\
 p\forces^{\mathrm{rec}}_{M,P}\varphi\to\psi
 &\Longleftrightarrow
   \forall q\leq_P p\ (q\forces^{\mathrm{rec}}_{M,P}\varphi
                  \Rightarrow q\forces^{\mathrm{rec}}_{M,P}\psi),\\
 p\forces^{\mathrm{rec}}_{M,P}\forall x\,\varphi(x,\bar\tau)
 &\Longleftrightarrow
   \forall\sigma\in M\ (\mathsf{Name}_P(\sigma)
        \Rightarrow p\forces^{\mathrm{rec}}_{M,P}\varphi(\sigma,\bar\tau)).
\end{aligned}
\]
\end{hbtheorem}

\subsection{Truth without a circular semantic definition}

\begin{hblemma}{thm:p06-ftr-atomic-truth}{proved=annex:p06.ftr.atomic-truth}{The atomic truth lemma, including equality}
For ground names \(\sigma,\tau\) and the supplied generic filter \(G\),
\[
\begin{aligned}
 \sigma^G\in\tau^G
 &\Longleftrightarrow\exists p\in G\
       (p\forces^{\mathrm{rec}}_{M,P}\sigma\in\tau),\\
 \sigma^G=\tau^G
 &\Longleftrightarrow\exists p\in G\
       (p\forces^{\mathrm{rec}}_{M,P}\sigma=\tau).
\end{aligned}
\]
The same equivalence holds for the negations of these atoms. The equality
direction uses a dense set of equality conditions and explicit disagreement
witnesses; it does not intersect infinitely many separately chosen conditions
from \(G\).
\end{hblemma}

\begin{hbtheorem}{thm:p06-ftr-truth}{proved=annex:p06.ftr.truth}{Forcing soundness and the truth lemma}
For every fixed standard formula and ground-name tuple,
\[
 (M[G],\in)\models\varphi(\bar\tau^G)
 \quad\Longleftrightarrow\quad
 \exists p\in G\ (p\forces^{\mathrm{rec}}_{M,P}\varphi(\bar\tau)).
\]
Moreover each \(p\in G\) forcing the displayed formula makes it true in
\((M[G],\in)\). The proof is an induction on the actual formula, with
all primitive and expanded constructors accounted for. It does not assume
in advance that \(M[G]\) satisfies \(\mathsf{ZFC}\).
\end{hbtheorem}

\begin{hbtheorem}{thm:p06-ftr-zfc}{proved=annex:p06.ftr.zfc}{Every ZFC axiom transfers to the generic extension}
Under \cref{def:p06-ftr-ground}, \((M[G],\in)\models\mathsf{ZFC}\).
This includes every Separation and Replacement instance with parameters
from \(M[G]\). Power Set is the power set computed in this structure,
not the ambient collection of every subset. The proofs construct exact
pair and union names, a separating subname, a Collection-bounded set of
image candidates, a bounded family of subset names, and a ground-indexed
surjection for Choice. Neither the assertion that a candidate is a name
nor the phrase ``rank recursion'' replaces its set-existence proof.
\end{hbtheorem}

\begin{hbcorollary}{thm:p06-ftr-semantic}{proved=annex:p06.ftr.semantic}{The externally countable semantic comparison}
If, in addition, \(M\) is externally countable, then for every \(p\in P\)
and fixed formula,
\[
 p\forces^{\mathrm{rec}}_{M,P}\varphi(\bar\tau)
 \quad\Longleftrightarrow\quad
 \text{every ambient }M\text{-generic }H\ni p
       \text{ satisfies }(M[H],\in)\models\varphi(\bar\tau^H).
\]
This is the earlier semantic relation at its stated countable-transitive
hypothesis. Countability supplies a generic through any countercondition;
it is not a premise of recursive definability or of the truth lemma for a
generic already supplied.
\end{hbcorollary}

\begin{hbexample}{def:p06-ftr-generic-boundary}{Why a filter is not enough}
In Cohen forcing \(P=2^{<\omega}\), stronger toward end-extension, put
\(D=\{s:\operatorname{length}(s)\geq1\}\) and
\(\dot G=\{\langle\check s,s\rangle:s\in P\}\).
Below every condition there is an \(s\in D\), and every strengthening of
\(s\) forces \(\check s\in\dot G\) and
\(\check s\in\check D\). Thus the existential clause forces
\(\exists x(x\in\dot G\land x\in\check D)\) at the greatest condition.
The nongeneric filter \(H=\{1_P\}\) nevertheless has
\(\dot G^H=\{1_P\}\), so that sentence is false in \(M[H]\).
The calculation uses the canonical-name atomic computation established in
\cref{annex:p06.ftr.atomic-truth}; it isolates the missing dense-set meeting,
not a failure of name evaluation. The full construction above applies to
arbitrary set forcing, not just to this example.
\end{hbexample}

\begin{proofboundary}{def:p06-ftr-source-boundary}{The remaining source and consistency boundaries}
These are local modern proofs at the stated transitive-ground/generic-filter
hypotheses. They do not transcribe Cohen's original relation, identify
Unger's auxiliary non-density-closed relation with this one, or discharge
the full separately allocated Jech and Kunen source presentations.
They also do not construct a generic over an arbitrary transitive ground,
prove the existence of a transitive model of \(\mathsf{ZFC}\), or derive
the bare arithmetized relative-consistency theorem from such a model.
The finite-fragment, syntactic and Boolean-valued discharges retain their
own obligations. Class forcing, iterations and forcing-axiom consistency
constructions are not consequences of this set-forcing entry.
\end{proofboundary}
